% GENERATED FILE. Edit canonical lessons, then run npm run generate:books.
% canonical-source-hash: fnv1a64:28d545e7741dc104
% canonical-lessons: RU-C227-shtraf, RU-C227-chelovek, RU-C227-muzhchina, RU-C227-zhenshchina, RU-C227-malchik

\chapter{Fine, Person, Man, Woman, Boy}
\label{ch:ru-fine-person-man-woman-boy}
\hlchaptermodality{227}

\begin{chapteropening}
\textbf{By the end of this chapter:} \emph{I can say a fine, a person, a man, a woman and a boy.}

Complete the last lesson of chapter 227: I can say a fine, a person, a man, a woman and a boy.
\end{chapteropening}

\section[\texorpdfstring{\ru{штраф} (shtraf)}{штраф (shtraf)}]{\ru{штраф} (shtraf) — a fine}

\label{lesson:RU-C227-shtraf}



\begin{quote}
Before the new one: say the Russian for a novel, then the Russian for poems, poetry.
\end{quote}

\subsection*{You'll want to know: \ru{штраф}}
\textbf{\ru{штраф}} (\emph{shtraf}) — \textquotedblleft{}a fine\textquotedblright{}.

\begin{grammarlens}[title={What you've built}]
One more word for this chapter.
\end{grammarlens}

\subsection*{Your turn}
\begin{itemize}
\raggedright
  \item \emph{Say it:} \emph{shtraf}
  \item \emph{Say it:} \emph{shtraf}, once more
  \item \emph{Say it:} say \emph{shtraf}
  \item \emph{Recall:} say the Russian for a novel, then the Russian for poems, poetry, then say \emph{shtraf} again
\end{itemize}

\begin{tcolorbox}[breakable,colback=teal!4,colframe=teal!35!black,title={Before you move on}]
What does \ru{штраф} mean? (\textquotedblleft{}A fine\textquotedblright{}.) Say it once more.
\end{tcolorbox}
\section[\texorpdfstring{\ru{человек} (chelovék)}{человек (chelovék)}]{\ru{человек} (chelovék) — a person}

\label{lesson:RU-C227-chelovek}



\begin{quote}
Before the new one: say the Russian for poems, poetry, then the Russian for a fine.
\end{quote}

\subsection*{You'll want to know: \ru{человек}}
\textbf{\ru{человек}} (\emph{chelovék}) — \textquotedblleft{}a person\textquotedblright{}.

\begin{grammarlens}[title={What you've built}]
One more word for this chapter.
\end{grammarlens}

\subsection*{Your turn}
\begin{itemize}
\raggedright
  \item \emph{Say it:} \emph{chelovék}
  \item \emph{Say it:} \emph{chelovék}, once more
  \item \emph{Say it:} say \emph{chelovék}
  \item \emph{Recall:} say the Russian for poems, poetry, then the Russian for a fine, then say \emph{chelovék} again
\end{itemize}

\begin{tcolorbox}[breakable,colback=teal!4,colframe=teal!35!black,title={Before you move on}]
What does \ru{человек} mean? (\textquotedblleft{}A person\textquotedblright{}.) Say it once more.
\end{tcolorbox}
\section[\texorpdfstring{\ru{мужчина} (muzhchína)}{мужчина (muzhchína)}]{\ru{мужчина} (muzhchína) — a man}

\label{lesson:RU-C227-muzhchina}



\begin{quote}
Before the new one: say the Russian for a fine, then the Russian for a person.
\end{quote}

\subsection*{You'll want to know: \ru{мужчина}}
\textbf{\ru{мужчина}} (\emph{muzhchína}) — \textquotedblleft{}a man\textquotedblright{}.

\begin{grammarlens}[title={What you've built}]
One more word for this chapter.
\end{grammarlens}

\subsection*{Your turn}
\begin{itemize}
\raggedright
  \item \emph{Say it:} \emph{muzhchína}
  \item \emph{Say it:} \emph{muzhchína}, once more
  \item \emph{Say it:} say \emph{muzhchína}
  \item \emph{Recall:} say the Russian for a fine, then the Russian for a person, then say \emph{muzhchína} again
\end{itemize}

\begin{tcolorbox}[breakable,colback=teal!4,colframe=teal!35!black,title={Before you move on}]
What does \ru{мужчина} mean? (\textquotedblleft{}A man\textquotedblright{}.) Say it once more.
\end{tcolorbox}
\section[\texorpdfstring{\ru{женщина} (zhénshchina)}{женщина (zhénshchina)}]{\ru{женщина} (zhénshchina) — a woman}

\label{lesson:RU-C227-zhenshchina}



\begin{quote}
Before the new one: say the Russian for a person, then the Russian for a man.
\end{quote}

\subsection*{You'll want to know: \ru{женщина}}
\textbf{\ru{женщина}} (\emph{zhénshchina}) — \textquotedblleft{}a woman\textquotedblright{}.

\begin{grammarlens}[title={What you've built}]
One more word for this chapter.
\end{grammarlens}

\subsection*{Your turn}
\begin{itemize}
\raggedright
  \item \emph{Say it:} \emph{zhénshchina}
  \item \emph{Say it:} \emph{zhénshchina}, once more
  \item \emph{Say it:} say \emph{zhénshchina}
  \item \emph{Recall:} say the Russian for a person, then the Russian for a man, then say \emph{zhénshchina} again
\end{itemize}

\begin{tcolorbox}[breakable,colback=teal!4,colframe=teal!35!black,title={Before you move on}]
What does \ru{женщина} mean? (\textquotedblleft{}A woman\textquotedblright{}.) Say it once more.
\end{tcolorbox}
\section[\texorpdfstring{\ru{мальчик} (málchik)}{мальчик (málchik)}]{\ru{мальчик} (málchik) — a boy}

\label{lesson:RU-C227-malchik}



\begin{quote}
Before the new one: say the Russian for a man, then the Russian for a woman.
\end{quote}

\subsection*{You'll want to know: \ru{мальчик}}
\textbf{\ru{мальчик}} (\emph{málchik}) — \textquotedblleft{}a boy\textquotedblright{}.

\begin{grammarlens}[title={What you've built}]
One more word for this chapter.
\end{grammarlens}

\subsection*{Your turn}
\begin{itemize}
\raggedright
  \item \emph{Say it:} \emph{málchik}
  \item \emph{Say it:} \emph{málchik}, once more
  \item \emph{Say it:} say \emph{málchik}
  \item \emph{Recall:} say the Russian for a man, then the Russian for a woman, then say \emph{málchik} again
\end{itemize}

\begin{tcolorbox}[breakable,colback=teal!4,colframe=teal!35!black,title={Before you move on}]
What does \ru{мальчик} mean? (\textquotedblleft{}A boy\textquotedblright{}.) Say it once more.
\end{tcolorbox}
